(1) 由每束 \(\partial_t\delta n_j+v_j\partial_x\delta n_j+n_b\partial_x\delta u_j=0\) 和 \((\partial_t+v_j\partial_x)\delta u_j=-(e/m)\delta E\)，代入 \(ik\delta E=-e(\delta n_++\delta n_-)/\varepsilon_0\) 得 \(\boxed{1-\omega_b^2/(\omega-kv_0)^2-\omega_b^2/(\omega+kv_0)^2=0}\)。

(2) 设 \(a=kv_0\)，清分母得到 \((\omega^2-a^2)^2=2\omega_b^2(\omega^2+a^2)\)。故 \(\omega_\pm^2=a^2+\omega_b^2\pm\omega_b\sqrt{\omega_b^2+4a^2}\)。下支为负当且仅当 \(0<a^2<2\omega_b^2\)，对应纯指数增长和衰减，阈值 \(kv_0=\sqrt2\omega_b\)。

(3) 不稳定范围增长率 \(\gamma=\sqrt{\omega_b\sqrt{\omega_b^2+4a^2}-a^2-\omega_b^2}\)。将 \(a^2\) 作为变量求最大，\(d\gamma^2/d(a^2)=2\omega_b/\sqrt{\omega_b^2+4a^2}-1=0\)，故 \(a^2=3\omega_b^2/4\)、\(\gamma_{\max}=\omega_b/2\)。\(a\to0\) 时下支 \(\omega_-^2\to0^-\)，这对应冷流的中性极限。

(4) 失稳电场的 \(e^{\gamma t}\) 增长无法从正定无源耗散模型产生，它以反向电子流的定向动能为自由能来源。对热分布必须用Vlasov介电函数 \(\varepsilon(\omega,k)\) 并处理速度空间共振与Landau阻尼；不能直接将零温有理函数的阈值用于高温流束。
